% 第 16 题转盘图：按固定源 PDF 矢量路径换算（1 pt → 0.03515 cm）。
% 原点取图组左下（源页 (1079.39, 554.08) pt），y 向上。圆半径 34.01，圆心即两直径交点；
% 两条 ±45° 直径把圆四等分；指针线在圆正上方、实心箭头向下（尖端距圆顶 3.40 pt 不接触）；
% 四扇区数字上 1 / 左 3 / 右 4 / 下 2；圆心实心点（源内 U+2022 字符，照录为转轴记号）。
\begin{tikzpicture}[line cap=round,line join=round,font=\small,x=0.03515cm,y=0.03515cm]
% 圆与两条垂直直径（均在圆内，端点落于圆周）
\draw (34.00,34.01) circle (34.01);
\draw (9.96,9.97)--(58.05,58.05);
\draw (58.05,9.97)--(9.96,58.05);
% 指针（竖线 + 向下实心箭头）
\draw (34.00,81.62)--(34.00,75.00);
\fill (32.51,75.40)--(35.50,75.40)--(34.005,71.41)--cycle;
% 圆心实心点（转轴记号）
\fill (34.43,34.41) circle (1.45);
% 标注（照源字形框中心，显式坐标）
\node[anchor=center,inner sep=0] at (34.43,88.90) {指针};
\node[anchor=center,inner sep=0] at (34.43,50.17) {\(1\)};
\node[anchor=center,inner sep=0] at (17.42,33.16) {\(3\)};
\node[anchor=center,inner sep=0] at (51.44,33.16) {\(4\)};
\node[anchor=center,inner sep=0] at (34.43,16.15) {\(2\)};
\end{tikzpicture}
